% Continuation frame 1: original text and reveals retained.
\begin{frame}[plain]
\PartStart{Light}{\PartSectionNumber\ / COLLABORATIVE FILTERING}{Collaborative filtering}{2305033 - collaborative}
\begin{PartBody}
\centering
\begin{tikzpicture}
\node[user] (u1) at (2,3.8) {U1}; \node[user] (u2) at (2,2.75) {U2}; \node[user] (u3) at (2,1.7) {U3}; \node[user] (u4) at (2,.65) {U4};
\node[item] (m1) at (7,4.15) {M1}; \node[item] (m2) at (7,3.1) {M2}; \node[item] (m3) at (7,2.05) {M3}; \node[item] (m4) at (7,1.0) {M4};
\draw[greyarr,visible on=<1->] (u1)--(m1); \draw[greyarr,visible on=<1->] (u1)--(m2); \draw[greyarr,visible on=<1->] (u2)--(m3); \draw[greyarr,visible on=<1->] (u2)--(m1); \draw[greyarr,visible on=<1->] (u2)--(m2); \draw[greyarr,visible on=<1->] (u3)--(m3); \draw[greyarr,visible on=<1->] (u4)--(m4);
\draw[dashedarr,visible on=<2->] (u1)--node[above,sloped,font=\scriptsize,text=rose]{predict}(m3);
\node[boxteal,text width=4.1cm,visible on=<3->] at (11.2,2.4) {\textbfPro{Collaborative signal}\\If users with similar histories liked $M_3$, estimate that $U_1$ may also like $M_3$.};
\end{tikzpicture}
\end{PartBody}
\end{frame}

% Continuation frame 2: original text and reveals retained.
\begin{frame}[plain]
\PartStart{Light}{\PartSectionNumber\ / TWO FAMILIES}{Collaborative filtering has two main families}{2305033 - two families}
\begin{PartBody}
\centering
\begin{tikzpicture}[outer sep=0pt]
\node[boxteal,text width=5.8cm,visible on=<1->] (r) at (6.7,5.3) {\textbfPro{Shared evidence}\\sparse user--item matrix $R$};
\node[boxblue,text width=5.4cm,minimum height=3.0cm,visible on=<2->] (mem) at (3.25,2.35) {\textbfPro{Memory-based}\\[3pt]Use $R$ directly\\Find similar users or items\\Predict from neighbours\\[3pt]\micro{User-KNN \textbullet\ Item-KNN}};
\node[boxviolet,text width=5.4cm,minimum height=3.0cm,visible on=<3->] (model) at (10.15,2.35) {\textbfPro{Model-based}\\[3pt]Learn parameters from $R$\\Compress users and items\\Predict from learned vectors\\[3pt]\micro{Matrix factorization \textbullet\ SVD}};
\draw[arr,visible on=<2->] ($(r.south west)!0.2!(r.south east)$)--(mem.north);
\draw[arr,visible on=<3->] ($(r.south west)!0.8!(r.south east)$)--(model.north);
\node[boxamber,text width=12.6cm,visible on=<4->] at (6.7,0.0) {\textbfPro{Same evidence:} neighbours now, or learned structure first.};
\end{tikzpicture}
\end{PartBody}
\end{frame}

% Continuation frame 3: original text and reveals retained.
\begin{frame}[plain]
\PartStart{Light}{\PartSectionNumber\ / MEMORY-BASED METHODS}{Memory-based methods}{2305033 - memory based}
\begin{PartBody}
\centering
\begin{tikzpicture}
\node[user] (target) at (6,2.7) {U};
\foreach \x/\y/\n in {4.3/3.9/A,4.8/1.55/B,7.25/3.6/C,7.8/1.6/D,5.9/4.2/E}{\node[user] (u\n) at (\x,\y) {\n};}
\draw[dashedarr,visible on=<2->] (target)--(uA); \draw[dashedarr,visible on=<2->] (target)--(uB);
\draw[greyarr,visible on=<3->] (uA)--(8.15,2.56); \draw[greyarr,visible on=<3->] (uB)--(8.15,2.56);
\node[item,visible on=<3->] (m) at (8.15,2.56) {M};
\node[boxblue,text width=3cm,visible on=<1->] at (1.6,2.55) {\textbfPro{User-based}\\find users similar to target user};
\node[boxteal,text width=3cm,visible on=<1->] at (10.8,3.25) {\textbfPro{Item-based}\\find items similar to target item};
\node[boxteal,text=ink,line width=.7pt,anchor=north,text width=12cm,visible on=<4->] at (6.5,.9) {Prediction comes from observed neighbours, without first learning a latent representation.};
\end{tikzpicture}
\end{PartBody}
\end{frame}

% Continuation frame 4: original text and reveals retained.
\begin{frame}[plain]
\PartStart{Light}{\PartSectionNumber\ / LIMITATIONS}{Why memory-based methods struggle at scale}{2305033 - limitations}
\begin{PartBody}
\centering
\begin{tikzpicture}[node distance=.75cm]
\node[boxrose,text width=2.9cm,visible on=<1->] (s) {\textbfPro{Sparsity}\\two users may overlap on very few items};
\node[boxamber,text width=3.2cm,right=of s,visible on=<2->] (scale) {\textbfPro{Scalability}\\similarity over millions of users/items is expensive};
\node[boxviolet,text width=3.2cm,right=of scale,visible on=<3->] (partial) {\textbfPro{Preference is partial}\\a click does not always mean preference};
\node[boxteal,text width=4.6cm,below=1.1cm of scale,visible on=<4->] (sol) {\textbfPro{Model-based direction}\\Learn a compact model once, then reuse it for prediction.};
\draw[arr,visible on=<4->] (s)--(sol); \draw[arr,visible on=<4->] (scale)--(sol); \draw[arr,visible on=<4->] (partial)--(sol);
\end{tikzpicture}
\end{PartBody}
\end{frame}

% Continuation frame 5: original text and reveals retained.
% \begin{frame}[plain]
% \PartStart{Light}{\PartSectionNumber\ / MATRIX FORMULATION}{Model-based matrix formulation}{2305033 - model formulation}
% \begin{PartBody}
% \centering
% \begin{tikzpicture}
% \node[boxrose,inner sep=6pt,visible on=<1->] (r) at (2.4,2.75) {\begin{tabular}{ccccc}5&?&4&?&1\\?&4&?&5&?\\4&?&?&4&2\\?&1&2&?&4\\5&?&?&?&4\end{tabular}};
% \node[font=\LARGE\bfseries,text=blue,visible on=<2->] at (5.0,2.75) {$\approx$};
% \node[boxblue,inner sep=6pt,visible on=<2->] (p) at (7.1,2.75) {\begin{tabular}{ccc}.6&.1&.8\\.7&.2&.7\\.1&.9&.9\\.9&.1&.1\\.2&.7&.8\end{tabular}};
% \node[font=\LARGE\bfseries,text=blue,visible on=<2->] at (8.9,2.75) {$\times$};
% \node[boxteal,inner sep=6pt,visible on=<2->] (q) at (11.4,2.75) {\begin{tabular}{ccccc}.9&.5&.7&.8&.2\\.1&.8&.2&.9&.3\\.5&.1&.8&.9&.3\end{tabular}};
% \node[boxteal,text=ink,line width=.7pt,anchor=north,text width=11.8cm,visible on=<3->] at (7,1.15) {Learn $P$ and $Q$ so that $p_u^\top q_i$ matches known values in $R$ and predicts the missing ones.};
% \end{tikzpicture}
% \end{PartBody}
% \end{frame}

% Continuation frame 6: original text and reveals retained.
\begin{frame}[plain]
\PartStart{Light}{\PartSectionNumber\ / LATENT FEATURES}{Model-based latent features}{2305033 - model representation}
\begin{PartBody}
\centering
\begin{tikzpicture}
\node[boxrose,text width=2.0cm,visible on=<1->] (r) at (1.6,3.1) {\textbfPro{Large sparse matrix}\\$m\times n$\\mostly unknown};
\node[boxblue,text width=2.0cm,visible on=<2->] (p) at (5.1,3.1) {\textbfPro{User matrix}\\$P$\\$m\times k$};
\node[boxteal,text width=2.0cm,visible on=<2->] (q) at (8.4,3.1) {\textbfPro{Item matrix}\\$Q^\top$\\$k\times n$};
\draw[arr,visible on=<2->] (r)--(p); \draw[arr,visible on=<2->] (p)--(q);
\node[boxamber,text width=3.2cm,visible on=<3->] at (11.8,2.7) {{\editorial\fontsize{14}{16}\selectfont Why possible?}\\User preferences are not random; they contain reusable patterns.};
\node[boxteal,text=ink,line width=.7pt,anchor=north,text width=10.8cm,visible on=<4->] at (6.5,.55) {The model may discover axes such as ``mainstream vs niche'' or ``light comedy vs dark drama'' without human labels.};
\end{tikzpicture}
\end{PartBody}
\end{frame}

% Transition from the methodology to the implementation algorithms.
\begin{frame}[plain]
\PartStart{Light}{\PartSectionNumber\ / FROM METHOD TO ALGORITHM}{How do these methods work in practice?}{2305033 $\rightarrow$ 2305052}
\begin{PartBody}
\centering
\begin{tikzpicture}
\node[boxteal,text width=10.8cm,minimum height=1.25cm,visible on=<1->] (method) at (6.5,4.05)
  {{\editorial\fontsize{15}{17}\selectfont We have seen the methodology}\\[3pt]
   Collaborative evidence can be used directly or compressed into latent factors.};
% \node[font=\fontsize{15}{17}\selectfont\bfseries,text=accent,text width=13.2cm,align=center,visible on=<2->] at (6.5,2.72)
%   {But how do these ideas become working algorithms?};
\node[boxblue,text width=4.1cm,minimum height=1.25cm,visible on=<2->] (knn) at (2.85,1.35)
  {\textbfPro{KNN}\\Memory-based prediction};
\node[boxamber,text width=4.1cm,minimum height=1.25cm,visible on=<3->] (svd) at (9.85,1.50)
  {\textbfPro{SVD}\\Model-based prediction};
\draw[arr,-{Latex[length=2.8mm]},visible on=<2->] (6.35,2.97) -- (knn.north east);
\draw[arr,-{Latex[length=2.8mm]},visible on=<3->] (6.35,2.97) -- (svd.north west);
\end{tikzpicture}
\end{PartBody}
\end{frame}

% Continuation: the next presenter begins with KNN.
\begin{frame}[plain]
\PartStart{Light}{\PartSectionNumber\ / NEAREST NEIGHBOURS}{KNN predicts from nearest neighbours}{2305052 - KNN}
\begin{PartBody}
\begin{columns}[T,onlytextwidth]
\column{.45\textwidth}
\centering
\begin{tikzpicture}[scale=.93]
\node[user] (t) at (2.3,2.6) {T};
\node[user] (a) at (1.2,3.6) {A}; \node[user] (b) at (3.6,3.35) {B}; \node[user] (c) at (1.0,1.7) {C}; \node[user] (d) at (3.7,1.35) {D};
\draw[dashedarr,visible on=<2->] (t)--(a); \draw[dashedarr,visible on=<2->] (t)--(b); \draw[dashedarr,visible on=<2->] (t)--(c);
\node[boxviolet,visible on=<2->] at (2.35,.35) {$k=3$ neighbors};
\end{tikzpicture}
\column{.55\textwidth}
\visible<3->{\callout[draw=violet!55,fill=paleViolet,align=center]{6.5cm}{\textbfPro{Weighted prediction}\\[3pt]$\displaystyle \hat r_{ui}=\bar r_u+\frac{\sum_{v\in N(u)}s(u,v)(r_{vi}-\bar r_v)}{\sum_{v\in N(u)} |s(u,v)|}$}}
\vspace{.15cm}
\visible<4->{\small The neighborhood provides local evidence, usually measured by cosine or Pearson similarity.}
\end{columns}
\end{PartBody}
\end{frame}

% The cosine-similarity recap was removed because it was already introduced earlier.
% The next presenter continues with the model-based algorithm.
\begin{frame}[plain]
\PartStart{Light}{\PartSectionNumber\ / LATENT FACTORS}{SVD learns latent factors}{2305052 - SVD}
\begin{PartBody}
\centering
\begin{tikzpicture}
\node[boxrose,text width=2.0cm,visible on=<1->] at (0.3,3.15) {\textbfPro{Observed}\\$R$\\$m\times n$};
\node[font=\LARGE\bfseries,text=blue,visible on=<2->] at (2.15,3.15) {$\approx$};
\node[boxblue,text width=1.5cm,visible on=<2->] at (3.75,3.15) {$U_k$\\$m\times k$};
\node[boxamber,text width=1.5cm,visible on=<2->] at (5.98,3.15) {$\Sigma_k$\\$k\times k$};
\node[boxteal,text width=1.5cm,visible on=<2->] at (8.15,3.15) {$V_k^\top$\\$k\times n$};
\node[boxviolet,text width=2.15cm,visible on=<3->] at (10.79,3.15) {\textbfPro{Prediction}\\$\hat r_{ui}=p_u^\top q_i$};
\node[boxviolet,text=ink,line width=.7pt,anchor=north,text width=10.4cm,visible on=<4->] at (6.4,1.48) {$\displaystyle \min_{P,Q}\sum_{(u,i)\in\Omega}(r_{ui}-p_u^\top q_i)^2+\lambda(\lVert p_u\rVert^2+\lVert q_i\rVert^2)$};

\end{tikzpicture}
\end{PartBody}
\end{frame}

% Continuation frame 10: original text and reveals retained.
\begin{frame}[plain]
\PartStart{Light}{\PartSectionNumber\ / RECOMMENDER WORKFLOW}{End-to-end recommender workflow}{2305052 - workflow}
\begin{PartBody}
\centering
\begin{tikzpicture}
\node[boxblue,visible on=<1->] (a) at (0.5,3.25) {Sparse data};
\node[boxteal,visible on=<2->] (b) at (2.98,3.25) {Train\\\micro{SGD / ALS}};
\node[boxamber,visible on=<3->] (c) at (5.69,3.25) {Latent vectors\\$P,Q$};
\node[boxviolet,visible on=<4->] (d) at (8.39,3.25) {Predict\\$\hat r_{ui}$};
\node[boxrose,visible on=<5->] (e) at (10.33,3.25) {Rank\\top-$N$};
\draw[arr,visible on=<2->] (a)--(b); \draw[arr,visible on=<3->] (b)--(c); \draw[arr,visible on=<4->] (c)--(d); \draw[arr,visible on=<5->] (d)--(e);
\node[boxblue,text width=2.75cm,visible on=<6->] at (0.75,1.25) {\textbfPro{KNN}\\Explainable, weaker at scale};
\node[boxteal,text width=2.75cm,visible on=<7->] at (5.12,1.25) {\textbfPro{MF / SVD}\\Compressed latent factors};
\node[boxviolet,text width=3.05cm,visible on=<8->] at (9.5,1.25) {\textbfPro{Modern systems}\\Hybrid: content + collaborative};
\end{tikzpicture}
\end{PartBody}
\end{frame}

% Continuation frame 11: original text and reveals retained.
\begin{frame}[plain]
\PartStart{Light}{\PartSectionNumber\ / TAKEAWAYS}{Takeaways}{Conclusion}
\begin{PartBody}
\centering
\begin{tikzpicture}
\node[box,draw=forest,fill=softgreen,text=charcoal,text width=2.65cm,visible on=<1->] (a) at (0.45,2.7) {\textbfPro{Sparse matrix}\\Observed only partially.};
\node[box,draw=forest,fill=softgreen,text=charcoal,text width=2.65cm,visible on=<2->] (b) at (4.10,2.7) {\textbfPro{Two families}\\Content + collaborative.};
\node[box,draw=forest,fill=softgreen,text=charcoal,text width=2.65cm,visible on=<3->] (c) at (7.65,2.7) {\textbfPro{Scaling idea}\\Compress into latent factors.};
\node[box,draw=forest,fill=softgreen,text=charcoal,text width=2.65cm,visible on=<4->] (d) at (11.17,2.7) {\textbfPro{Final output}\\Ranked personalized choices.};
\end{tikzpicture}
\vspace{.35cm}
\visible<5->{\callout[rounded corners=2.2mm,draw=goldink,fill=softgold,text=charcoal]{11.6cm}{\textbfPro{The system does not know exactly what we will like \\ it estimates from patterns shared across users, items and context.}\par\vspace{2pt}}}
\end{PartBody}
\end{frame}
